\section{Tartışma}
\label{sec:discussion}

\subsection{Yinelemeli Isınma Gecikmesi}
\label{sec:disc:warmup}

Özgün \fastgrnn{} makalesinde~\cite{kusupati2018fastgrnn} ele
alınmayan önemli bir nokta şudur: sınıf etiketinin kararlı hâle
gelmesi için, kullanılan donanımdan bağımsız olarak, yaklaşık iki
saniyelik ($\sim$100 örnek, 50~Hz) gizli durum geçmişi gerekir.
Şekil~\ref{fig:warmup}, tek bir \textsc{standing} penceresi boyunca
üretilen etiketi ve $h_0$ bileşenini izler. Model, $t{=}100$
civarında doğru sınıfa kilitlenmeden önce \textsc{walking} ve
\textsc{upstairs} tahminlerinden geçer. Bu davranış iki MCU'dan
birine özgü değildir; Arduino ve MSP430 aynı yörüngeleri üretir
(Bölüm~\ref{sec:results:crossplatform}).

Düşme algılama, jest tanıma ya da duruş koçluğu gibi kullanıcıya
dönük uygulamalarda bu bulgu önemlidir. Aktivite \emph{geçişleri},
ancak gerçekleştikten yaklaşık 2~s sonra güvenle doğrulanabilir. Bu
kayda değer bir gecikmedir; cihaz-üstü RNN'ler değerlendirilirken
örnek-başına hızın yanında ayrıca raporlanmalıdır. Benzer
büyüklükteki başka kapılı yinelemeli hücrelerde de aynı etkinin
görülebileceğini düşünüyoruz. Bunu eşleşen parametre sayılarında
LSTM/GRU taban çizgileriyle sınamak doğal bir devam çalışmasıdır.

\begin{figure}[t]
\centering
\includegraphics[width=\linewidth]{figures/warmup_curve.pdf}
\caption{Tek bir 2.56~s pencerede (50~Hz, \textsc{standing} gerçek
etiketi) gizli durum bileşeni $h_0$ ve üretilen sınıf etiketi.
Arduino ve MSP430 aynı yörüngeleri üretir. Yaklaşık 2~s geçmiş
birikene kadar sınıf etiketi kararsızdır.}
\label{fig:warmup}
\end{figure}

\subsection{Yeniden Üretilebilirlik Olarak Platformlar Arası Determinizm}
\label{sec:disc:determinism}

Modern ML çıkarımı, donanımlar arasında çoğu zaman birebir yeniden
üretilebilir değildir. Üreticiye özgü BLAS uygulamaları, karma
duyarlıklı tensör çekirdekleri ve birleşmeli olmayan kayan nokta
indirgemeleri küçük farklar doğurabilir. Buna karşılık taşınabilir C
motorumuz, hem 8-bit AVR hem de 16-bit MSP430 hedeflerinde yazılımla
öykünülen kayan nokta kullanmasına rağmen platformlar arasında
\emph{bit-eşdeğer} tahminler üretir
(Bölüm~\ref{sec:results:crossplatform}). Bu düzeyde determinizm,
tıbbi nitelikli HAR gibi güvenlikle ilişkili gömülü uygulamalarda
korunması gereken bir özelliktir. Düzenleyici kurumlar, bir yazılım
değişikliğinin yalnızca yetkilendirildiği donanım üzerinde yeniden
doğrulanmasını isteyebilir. İki farklı ISA üzerinde bit-eşdeğer kalan
bir çıkarım yolu, bu kısıtı hafifletmek açısından değerlidir.

\subsection{Zaman Payı Analizi ve Saf-Q15 Çıkmazı}
\label{sec:disc:headroom}

\msp{}, örnek başına 20~ms bütçenin \%65'ini kullanır ve
$\sim$7~ms zaman payı bırakır (Bölüm~\ref{sec:results:streaming}).
Canlı MPU-6050 testlerinde bu pay, örnek başına bir
I\textsuperscript{2}C ivmeölçer okumasını, bir LED durum
güncellemesini ve bir UART günlük satırını rahatça karşılar.

Daha agresif bir \emph{saf-Q15} çıkarım yolu da denendi: ağırlıklar,
aktivasyonlar \emph{ve} ara çarpımların tamamı Q15 tamsayı
aritmetiğinde tutuldu. Bu yol pratikte başarısız oldu. Çok aşamalı
düşük ranklı çarpım sırasında tensör-başına ölçekler zincirlenerek
$\sim$$10^{-13}$ düzeyine kadar küçülüyor; bu değer int64 biriktirici
ve özel bir aşama-başına kaydırma çizelgesi olmadan Q15 çarpıcıda
temsil edilemiyor. Bölüm~\ref{sec:method:q15}'de tartışılan sabit
Q9.6 biçimi ölçek zincirleme sorununu hafifletir; ancak bu kez
ağırlık tarafında çözünürlük kaybı yaratır. Ağırlık büyüklükleri
tensörler arasında dört büyüklük mertebesine yayıldığı için, tek bir
Q9.6 ölçeği bazı tensörlerde kalibre edilmiş yola göre $8\times$'a
varan çözünürlük kaybına neden olur. Daha temiz çözüm,
nicelendirme farkında eğitimdir~\cite{hubara2017qnn}; çünkü ölçek
çizelgesi eğitim sonrasında zorlanarak değil, eğitim sırasında
öğrenilir. LUT tabanlı FP32 yol 50~Hz gerçek zamanlı kısıtı yeterli
zaman payıyla karşıladığı için bu çalışmada QAT yolunu sürdürmedik.
Bunu gelecek çalışma olarak ele alıyoruz
(Bölüm~\ref{sec:disc:future}).

\subsection{Sınırlılıklar}
\label{sec:disc:limitations}

İddialarımızın kapsamını dört sınırlılık belirler:

\begin{itemize}
\item \textbf{Veri kümesi.} \hapt{}, sabit biçimde bele takılmış bir
akıllı telefonla laboratuvar koşullarında kaydedilmiştir. Sensörün
vücut üzerindeki yer değiştirmesi, serbest yaşam verilerindeki dağılım
kayması ve özneler arası değişkenliğin HAR doğruluğunu düşürdüğü
bilinmektedir~\cite{wang2019harsurvey,demrozi2020harsurvey}; bunlar
bu çalışmada modellenmemiştir.

\item \textbf{Etiket kümesi.} Model altı temel \hapt{} aktivitesi
üzerinde eğitilmiş ve değerlendirilmiştir. Tam \hapt{} etiket kümesi
altı postüral geçiş daha içerir (ör.\ \textsc{stand-to-sit},
\textsc{lie-to-sit}); modelimiz bu geçişleri işlemez. Bu tercih,
özgün \fastgrnn{} HAR kıyaslamasıyla eşleşmek için yapılmıştır.

\item \textbf{DOWNSTAIRS sınıfı.} \textsc{downstairs}, boru hattı
boyunca performansı belirleyen en zor sınıftır
(Şekil~\ref{fig:perclass}). Bu makalede sınıfa özgü bir iyileştirme
denemesi yapılmamıştır. Bölüm~\ref{sec:disc:future} olası birkaç
müdahaleyi özetler.

\item \textbf{İstatistiksel anlamlılık testleri.} Çok tohumlu
değerlendirmede beş tohum üzerinde ortalama$\pm$standart sapma
raporlanmıştır; ancak eşleştirilmiş anlamlılık testleri ya da
parametrik olmayan bootstrap aralıkları verilmemiştir. Biçimsel
anlamlılık testi gelecek çalışmaya bırakılmıştır.
\end{itemize}

\subsection{Gelecek Yönelimler}
\label{sec:disc:future}

Bu çalışmanın ardından izlenebilecek beş doğal yön vardır.

\begin{itemize}
\item \textbf{Statik ve dinamik örüntüler için çift-rank ayrıştırma.}
Rank taramamızda (Bölüm~\ref{sec:results:lowrank}) $r_u{=}4$ dinamik
sınıflar (\textsc{walking}, \textsc{upstairs}, \textsc{downstairs})
için, $r_u{=}8$ ise statik sınıflar (\textsc{sitting},
\textsc{standing}) için daha elverişlidir. Düşük maliyetli bir
varyant olarak
$\mathbf{U}_\text{eff} = \mathrm{LowRank}(r{=}4) + \mathrm{diag}(\boldsymbol{\alpha})$
kullanılabilir. Bu yapı yalnızca $H$ ek parametre ekler; statik
DC-benzeri sinyalin diyagonal artık üzerinden, dinamik örüntünün ise
düşük rank kolu üzerinden geçmesine izin verir. Bu yaklaşımın,
statik sınıflarda bozulma yaratmadan $r_u{=}4$'ün dinamik sınıf
başarımını geri kazanmasını bekliyoruz.

\item \textbf{DOWNSTAIRS hedefli öznitelikler.} Standart UCI~HAR
boru hattı, yerçekimini vücut ivmesinden ayırmak için 0.3~Hz'de
Butterworth alçak geçiren süzgeç uygular. Özgün \fastgrnn{} HAR
kurulumuyla eşleşmek için bu adımı kullanmadık. Bu süzgeci yeniden
eklemek ya da 128 örnekli pencere üzerinde az sayıda FFT-bant
özniteliğiyle girdiyi genişletmek, \textsc{downstairs} sınıfındaki
farkı kapatmanın düşük maliyetli bir yolu olabilir.

\item \textbf{Saf-Q15 çıkarım için nicelendirme farkında eğitim.}
Bölüm~\ref{sec:disc:headroom}'da tartışıldığı gibi, ölçek
zincirleme nedeniyle eğitim-sonrası nicelendirme tam tamsayı boru
hattında başarısız olur. QAT~\cite{hubara2017qnn} bu sorunu
çözebilir. Saf-Q15 yolun \msp{} üzerinde ek 2--3$\times$ hızlanma
sağlaması muhtemeldir ve LUT tabanlı yaklaşımın doğal devamıdır.

\item \textbf{Geçiş durumları.} Tam \hapt{} etiket kümesindeki altı
postüral geçişi eklemek, yinelemeli durumun uzun ufuklu aktivite
tanıma ile kısa ufuklu geçiş algılama arasında nasıl paylaştırılacağı
sorusunu doğurur. İki farklı örnekleme temposunda değerlendirilen
çok ölçekli bir yinelemeli hücre, makul bir başlangıç noktasıdır.

\item \textbf{Çıkarım başına enerji.} Bu çalışma çıkarım-başına zaman
raporlar; ancak çıkarım-başına joule raporlamaz. Kart üzerinde
doğrudan akım ölçümü, gömülü uygulama değerlendirmesini tamamlar ve
Bölüm~\ref{sec:intro}'daki motivasyona geri bağlanır. Silikon arzı
ve enerji kısıtları birlikte düşünülmelidir; 30$\times$ çalışma süresi
hızlanması, uygulama düzeyinde ancak pil tüketimini de orantılı
biçimde azaltıyorsa anlamlıdır.
\end{itemize}
